\documentclass[11pt]{article}
\usepackage[margin=1in]{geometry}
\usepackage{amsmath,amssymb,amsthm}
\usepackage{booktabs}
\usepackage{listings}
\usepackage{xcolor}
\usepackage[numbers]{natbib}
\usepackage{hyperref}
\sloppy

\newtheorem{theorem}{Theorem}
\newtheorem{lemma}[theorem]{Lemma}
\theoremstyle{definition}
\newtheorem{proposal}[theorem]{Proposed statement}
\theoremstyle{remark}
\newtheorem{remark}[theorem]{Remark}

\lstdefinelanguage{lean}{
  morekeywords={theorem,def,example,import,namespace,end,open,by,intro,exact,obtain,refine,have,simp,only,at,fun,noncomputable},
  sensitive=true, morecomment=[l]{--}, morecomment=[s]{/-}{-/}, morestring=[b]"}
\lstset{language=lean, basicstyle=\ttfamily\footnotesize, keywordstyle=\color{blue!60!black}\bfseries,
  commentstyle=\color{green!40!black}, columns=fullflexible, breaklines=true, frame=single,
  extendedchars=true, inputencoding=utf8,
  literate={≤}{{$\le$}}1 {≥}{{$\ge$}}1 {→}{{$\to$}}1 {↔}{{$\leftrightarrow$}}1 {ℝ}{{$\mathbb{R}$}}1
           {ℕ}{{$\mathbb{N}$}}1 {⟨}{{$\langle$}}1 {⟩}{{$\rangle$}}1 {∃}{{$\exists$}}1 {∀}{{$\forall$}}1
           {∧}{{$\wedge$}}1 {·}{{$\cdot$}}1 {−}{{$-$}}1 {Δ}{{$\Delta$}}1 {μ}{{$\mu$}}1 {σ}{{$\sigma$}}1
           {‖}{{$\|$}}1 {∑}{{$\sum$}}1 {≠}{{$\ne$}}1 {∈}{{$\in$}}1 {←}{{$\leftarrow$}}1}

\title{Re-certifying the Viridis OS Decision Kernels:\\ Carbon Continuity, and Sealed Statement Proposals\\ for Restoration, Afforestation and Harmonization}
\author{Justin Hart\\ Viridis LLC}
\date{30 October 2026}

\begin{document}
\maketitle

\begin{abstract}
Viridis OS runs decision kernels whose formal cores were checked by tools that can no longer be re-run.
This paper documents the project that brings each kernel to a re-checkable standard: the kernel's formal core is restated as a
sealed, self-contained Lean~4 statement and re-proved under a pinned toolchain.
The first kernel, \emph{Carbon Continuity}, has a Lean proof of its sealed statement: a two-pool carbon portfolio can be kept
from shrinking exactly when $(1-a)(1-d)\le rp$.
This revision adds proposed sealed statements for the next three kernels.
\texttt{P2.Restoration} restates the forest-nucleation core.
The critical radius $r^\star=2a\sigma/(3\Delta\mu)$ is positive and is the unique maximiser of the free energy
$G(r)=-\Delta\mu r^3+a\sigma r^2$ on $r\ge 0$, and the barrier satisfies $G(r^\star)=\tfrac12\Delta\mu (r^\star)^3$.
\texttt{P2.Afforestation} restates two results: seeding efficiency scales as the cube of the ecological drive, and the seeding
alignment $\eta=\langle n,m\rangle^2/(\|n\|^2\|m\|^2)$ lies in $[0,1]$.
\texttt{P2.Harmonization} restates two properties of a scalar-price welfare theorem for $N$ quadratic stewards: market clearing,
and the absence of a duality gap.
Each proposal copies the canon definitions and hypotheses verbatim and records the hash of the canon file.
It also includes sanity examples with concrete numbers, among them negative controls.
We point out where canon theorem names promise more than the theorems prove.
We also point out which canon hypotheses are kept for faithfulness although the proofs do not need them.
Sealing the proposals is a human decision, and no proof of any proposal is claimed here.
\end{abstract}

\section{Introduction}
Viridis OS runs a set of decision kernels. Each one is a small piece of software that turns ecological
parameters into a decision. A kernel's rule is only as trustworthy as the mathematics behind it.
The archived canon~\cite{HartCanon2026} proved those rules in Lean, but the certificates came from tooling that has
since been retired. The project described here moves each kernel's assurance onto a re-checkable footing in two steps.
First, a \emph{sealed statement} fixes the exact proposition a kernel relies on.
Second, a Lean proof of that statement is checked under Lean v4.28.0 with a pinned Mathlib~\cite{Mathlib2020}.

The two steps are separate on purpose. Writing the statement is where most of the risk lies.
A statement can be correct and easy to prove and still be the wrong one. It can be weaker than what the kernel uses,
stronger than what the canon proved, or vacuous. The present revision is about that first step for three kernels:
Restoration, Afforestation and Harmonization.
The question is narrow. We ask what each kernel's rule \emph{mathematically} relies on, and how to write that down
self-contained and faithful to the canon theorem it replaces. We do not ask whether the models describe any real landscape.

\section{Background and related work}
\paragraph{Carbon Continuity.} One disturbance-recovery cycle maps $(L,D)\mapsto(aL+rD,\ pL+dD)$.
This is the nonnegative matrix $M=\begin{pmatrix}a&r\\p&d\end{pmatrix}$, and the threshold theorem is a
two-dimensional instance of the Collatz--Wielandt characterisation in Perron--Frobenius theory~\cite{BermanPlemmons1994}.

\paragraph{Restoration and Afforestation.} The canon models a planted cluster as a nucleus of the forest phase inside a
metastable degraded phase. It uses the free-energy form of classical nucleation theory~\cite{Kashchiev2000}:
a bulk gain $-\Delta\mu\, n$ competes with a surface cost $a\sigma\, n^{2/3}$.
In the cube-root variable $r=n^{1/3}$ this becomes the cubic $G(r)=-\Delta\mu r^3+a\sigma r^2$.
The textbook facts are a critical size at the barrier top and a barrier equal to half the bulk term at the critical size.
These are the mathematical content of \texttt{P2.Restoration}.
The Afforestation efficiency $e=C/n^\star$ inherits the cube from $n^\star=(r^\star)^3$.
The alignment score is a squared cosine, bounded by Cauchy--Schwarz.

\paragraph{Harmonization.} The scalar-price theorem is the separable quadratic case of Lagrangian duality with
one linear coupling constraint~\cite{BoydVandenberghe2004}.
At the clearing price, the decentralised best responses solve the centralised problem, and there is no duality gap.

We claim no novelty for any of this mathematics. The contribution is a set of faithful, machine-checkable targets tied
to deployed kernels.

\section{Carbon Continuity (P2-T001)}
The sealed module defines $\mathrm{stepLiving}(a,r,L,D)=aL+rD$ and $\mathrm{stepDurable}(d,p,L,D)=pL+dD$.
$\mathrm{Sustainable}(a,d,r,p)$ means that some $L,D>0$ satisfy $L\le aL+rD$ and $D\le pL+dD$.
\begin{theorem}[\texttt{P2.CarbonContinuity}]
For all real $a,d,r,p$ with $0\le a<1$, $0\le d<1$, $r>0$ and $p>0$:
$\mathrm{Sustainable}(a,d,r,p)\iff (1-a)(1-d)\le rp$.
\end{theorem}
\begin{proof}[Proof sketch]
Sufficiency uses the witness $(L,D)=(r,1-a)$. The living step is then an identity, and the durable step reduces to
$rp-(1-a)(1-d)\ge 0$. For necessity, rearrange the two inequalities to $(1-a)L\le rD$ and $(1-d)D\le pL$,
multiply them (both sides are positive), and cancel $LD>0$.
\end{proof}
The previous revision gave a Lean proof of this statement (file \texttt{proofs/P2.CarbonContinuity.lean} of run
2026-10-29), with axioms \texttt{propext}, \texttt{Classical.choice} and \texttt{Quot.sound}.
\begin{lstlisting}
theorem main : ViridisStatements.P2.CarbonContinuity.Statement := by
  intro a d r p _ ha1 _ hd1 hr _
  exact ⟨necessary a d r p ha1 hd1, sufficient a d r p ha1 hr⟩
\end{lstlisting}
We observe that the helper lemmas use only $a<1$, $d<1$ and $r>0$.

\section{Method for writing a sealed statement}
For each kernel we followed the same procedure.
\begin{enumerate}
\item \textbf{Identify the canon theorems.} We take the theorems named in the task brief and check the canon file's SHA-256
against the pinned manifest (canon commit \texttt{80d32904}).
\item \textbf{Copy definitions verbatim.} We copy only the definitions the theorems use, with the same argument order,
and drop the rest. Any added abbreviation must unfold to the canon expression exactly.
\item \textbf{Copy hypotheses verbatim.} Each conjunct of \texttt{Statement} keeps exactly the hypotheses of its canon theorem.
We drop no hypothesis, even one that is mathematically unnecessary, because dropping it would strengthen the canon statement.
We add none, because adding one would weaken it.
\item \textbf{State the conclusion as proved, not as named.} Where a canon docstring or name claims more than the Lean
conclusion, the statement follows the conclusion.
\item \textbf{Pin the meaning.} We add sanity \texttt{example}s with concrete numbers, closed by \texttt{norm\_num} or \texttt{simp}
(\texttt{decide} does not apply to $\mathbb{R}$). They include negative controls showing that the key hypotheses or factors cannot be dropped.
\item \textbf{Test non-triviality.} We compile each module and probe \texttt{Statement} with 28 one-shot tactics
after \texttt{unfold Statement}. The tactics include \texttt{rfl}, \texttt{decide}, \texttt{simp}, \texttt{norm\_num}, \texttt{aesop},
\texttt{linarith}, \texttt{nlinarith}, \texttt{positivity}, \texttt{grind}, \texttt{field\_simp}, \texttt{ring} and \texttt{polyrith}.
\item \textbf{Test feasibility.} Privately, we check that a port of the canon proof closes the proposed statement.
This guards against proposing a false target. It is not a claim of proof.
\end{enumerate}

\section{The three proposals}

\subsection{\texttt{P2.Restoration}: the critical-nucleus go/no-go}
Definitions: $r^\star(\Delta\mu,a,\sigma)=2a\sigma/(3\Delta\mu)$ and $G(\Delta\mu,a,\sigma,r)=-\Delta\mu r^3+a\sigma r^2$.
\begin{proposal}
(i) For all $\Delta\mu,a,\sigma>0$: $r^\star>0$; $G(r)\le G(r^\star)$ for all $r\ge0$; and $G(r)=G(r^\star)$ with $r\ge0$ implies $r=r^\star$.
(ii) For all $\Delta\mu>0$ and all $a,\sigma$: $G(r^\star)=\tfrac12\Delta\mu\,(r^\star)^3$.
\end{proposal}
Part (i) restates canon \texttt{critical\_nucleus}, and part (ii) restates \texttt{barrier\_half\_identity}
(file sha256 \texttt{667a826a\ldots ec28132}).
Both follow from the factorisation $G(r^\star)-G(r)=\Delta\mu\,(r-r^\star)^2\,(r+r^\star/2)$.
The kernel's go/no-go rule compares cluster size with $n^\star=(r^\star)^3$. Below $n^\star$ the cluster is on the uphill side of the barrier.
Negative controls: $G(-1)=5/2>G(r^\star)=1/2$ at $(\Delta\mu,a,\sigma)=(1,1,3/2)$, so the restriction $r\ge0$ is essential.
$G(r^\star)\ne\Delta\mu(r^\star)^3$, so the factor $\tfrac12$ is substantive.
$\Delta\mu=-1$ gives $r^\star<0$.
\begin{lstlisting}
def Statement : Prop :=
  (∀ Δμ a σ : ℝ, 0 < Δμ → 0 < a → 0 < σ →
    0 < rStar Δμ a σ ∧
    (∀ r : ℝ, 0 ≤ r → G Δμ a σ r ≤ G Δμ a σ (rStar Δμ a σ)) ∧
    (∀ r : ℝ, 0 ≤ r → G Δμ a σ r = G Δμ a σ (rStar Δμ a σ) → r = rStar Δμ a σ)) ∧
  (∀ Δμ a σ : ℝ, 0 < Δμ →
    G Δμ a σ (rStar Δμ a σ) = (1 / 2) * Δμ * (rStar Δμ a σ) ^ 3)
\end{lstlisting}

\subsection{\texttt{P2.Afforestation}: cubic efficiency and seeding alignment}
Definitions: $\mathrm{eff}(a,\sigma,\Delta\mu,C)=C\,(3\Delta\mu/(2a\sigma))^3$ and $\eta(n,m)=\langle n,m\rangle^2/(\|n\|^2\|m\|^2)$.
\begin{proposal}
(i) For all $a,\sigma,\Delta\mu,t>0$ and all $C$: $\mathrm{eff}(a,\sigma,t\Delta\mu,C)=t^3\,\mathrm{eff}(a,\sigma,\Delta\mu,C)$.
(ii) In every real inner-product space $E$ (universe-polymorphic, as in the canon), for $n,m\ne0$: $0\le\eta(n,m)\le1$.
\end{proposal}
These restate the canon theorems
\texttt{establishment\_efficiency\_cubic\_in\_drive\_over\_tension} and \texttt{seeding\_efficiency\_eq\_cos2\_theta}
 (file sha256 \texttt{3ebe33bf\ldots 9526f}).
\begin{remark}
The second canon theorem is named ``$\mathrm{eff}=\cos^2\Theta$''. Its docstring also mentions $\eta=1$ iff $n\parallel m$, and $1-\eta=\sin^2\Theta$.
Its Lean conclusion proves only the bound $0\le\eta\le1$, and the proposal restates only that bound.
If the kernel uses the equality case, that case needs its own statement.
\end{remark}
Sanity examples pin $\mathrm{eff}=C/n^\star$ ($27/8$ at unit parameters), the factor $8$ when the drive is doubled, and the factor $8$ when $\sigma$ is halved.
They also pin $\eta=1$ for parallel and $\eta=0$ for orthogonal vectors in $\mathbb{R}^2$.
Negative controls: doubling the drive does not double the efficiency, and the unnormalised $\langle2,3\rangle^2=36$ exceeds $1$.
\begin{lstlisting}
def Statement : Prop :=
  (∀ a σ Δμ C t : ℝ, 0 < a → 0 < σ → 0 < Δμ → 0 < t →
    eff a σ (t * Δμ) C = t ^ 3 * eff a σ Δμ C) ∧
  (∀ (E : Type*) [NormedAddCommGroup E] [InnerProductSpace ℝ E] (n m : E),
    n ≠ 0 → m ≠ 0 → 0 ≤ eta n m ∧ eta n m ≤ 1)
\end{lstlisting}

\subsection{\texttt{P2.Harmonization}: one price suffices}
Definitions: $\mathrm{br}(\lambda)_i=b_i-\lambda/a_i$, $J(x)=\sum_i \tfrac{a_i}{2}(x_i-b_i)^2$,
$\lambda^\star=(\sum_i b_i-X)/\sum_i a_i^{-1}$ and $x^\star_i=\mathrm{br}(\lambda^\star)_i$.
\begin{proposal}
For all $N>0$, all $a,b\in\mathbb{R}^N$ with every $a_i>0$, and all $X$:
(i) $\sum_i \mathrm{br}(\lambda^\star)_i=X$; and
(ii) every $y$ with $\sum_i y_i=X$ satisfies $J(x^\star)\le J(y)$.
\end{proposal}
These restate \texttt{ght\_market\_clears} and \texttt{ght\_strong\_duality\_decentralized\_eq\_centralized}
(file sha256 \texttt{cfdead65\ldots 110aec}).
On the canon instance $N=2$, $a=b=(1,1)$, $X=0$, the sanity examples give $\lambda^\star=1$, $x^\star=(0,0)\ne b$, $J(x^\star)=1$, and $J(1,-1)=2$.
Negative controls:
\begin{itemize}
\item The infeasible ideal $y=b$ has $J=0<1$, so the feasibility hypothesis is essential.
\item The price $\lambda=0$ does not clear.
\item With $N=0$ the aggregate is $0\ne X=1$, so $N>0$ is essential.
\end{itemize}
\begin{lstlisting}
def Statement : Prop :=
  ∀ (N : ℕ) (a b : Fin N → ℝ), 0 < N → (∀ i, 0 < a i) → ∀ X : ℝ,
    (∑ i, br a b (lamStar a b X) i) = X ∧
    (∀ y : Fin N → ℝ, (∑ i, y i) = X → J a b (xStar a b X) ≤ J a b y)
\end{lstlisting}

\section{Results}
\begin{table}[h]
\centering\small
\begin{tabular}{@{}p{0.17\textwidth}p{0.42\textwidth}lc@{}}
\toprule
Proposal & Canon theorems restated & Canon file sha256 (prefix) & Probes closed \\
\midrule
\texttt{P2.Restoration} & \texttt{critical\_nucleus}, \texttt{barrier\_half\_identity} & \texttt{667a826a} & 0/28 \\
\texttt{P2.Afforestation} & \texttt{establishment\_efficiency\_\ldots}, \texttt{seeding\_efficiency\_\ldots} & \texttt{3ebe33bf} & 0/28 \\
\texttt{P2.Harmonization} & \texttt{ght\_market\_clears}, \texttt{ght\_strong\_duality\_\ldots} & \texttt{cfdead65} & 0/28 \\
\bottomrule
\end{tabular}
\caption{Proposed sealed statements. All three modules compile under the pinned toolchain, and their sanity examples elaborate.
The probe column counts how many of our 28 one-shot tactics closed \texttt{Statement} locally. The verifier runs its own probes.}
\end{table}
We observe the following.
\begin{itemize}
\item For each proposal, a private port of the canon proof elaborated against a local build of the module.
It depends only on \texttt{propext}, \texttt{Classical.choice} and \texttt{Quot.sound}.
We do not present these as proofs of the proposals. The proofs belong to later tasks, against the sealed modules.
\item Several canon hypotheses are not needed mathematically but are kept for faithfulness:
\begin{itemize}
\item In Lean's convention $x/0=0$, the half-barrier identity also holds at $\Delta\mu=0$.
\item The cubic law is a ring identity, so it needs none of its positivity hypotheses.
\item The alignment bound holds trivially when $n=0$ or $m=0$, since then $\eta=0$.
\end{itemize}
\item One canon name over-promises: the ``$\cos^2\Theta$'' theorem proves only the bound $0\le\eta\le1$.
\end{itemize}

\section{What the statements settle and what they do not}
\begin{itemize}
\item \textbf{Settled, once sealed and proved:} the internal mathematics of each kernel's rule inside its stated model.
\item \textbf{Not settled:} any parameter value. This covers $\Delta\mu$, $\sigma$, $a$, $C$, the curvatures $a_i$, the ideals $b_i$ and the target $X$.
These come from field evidence and remain ASSUMED.
\item \textbf{Not settled:} whether the models are adequate. The models in question are classical nucleation for tree-cluster establishment and
quadratic dissipation for stewardship coordination. Nucleation rates, ecological temperature and stochastic establishment are outside the statements.
\item \textbf{Not settled:} code conformance. The kernels' runtime identifiers and TypeScript sources were not among this session's inputs.
\end{itemize}

\section{Limitations}
The mathematics is elementary, and its value lies in being faithful and re-checkable.
Faithfulness was judged by hand against the canon text. No tool checks that a proposal matches its canon theorem.
Our probe list is our own guess at the verifier's and may differ from it.
\texttt{P2.Afforestation} is universe-polymorphic, following the canon's \texttt{Type*}. We checked locally that a polymorphic
candidate \texttt{theorem main : Statement} and the verifier's statement-match template both elaborate.
A universe-monomorphic seal (\texttt{Type}) would be a mild weakening.
The task brief mentions a guidance file (\texttt{STATEMENT\_PROPOSAL.md}) that was not among the inputs.
We followed the brief and the format of the existing sealed module instead.

\section{Reproducibility}
Toolchain \texttt{leanprover/lean4:v4.28.0}, Mathlib commit \texttt{8f9d9cff6bd728b17a24e163c9402775d9e6a365}.
Canon repository commit \texttt{80d32904d751b45f415e23f872510cc2eddca3ff}. The canon file hashes are listed in the proposals' rationale files.
Each proposal compiles with \texttt{lean proposals/P2.}\textit{Name}\texttt{.lean}.
The sealed Carbon Continuity statement has sha256 \texttt{a5a6e8c6\ldots ffded}. To check its proof, run
\texttt{viridis-lean-check proofs/P2.CarbonContinuity.lean P2.CarbonContinuity}.
No automated prover jobs and no external data were used in either session.

\section{Changelog}
\begin{itemize}
\item 2026-10-29 (P2-T001): First version. Ported the canon threshold proof to the sealed statement
\texttt{P2.CarbonContinuity}, and wrote the scope statement.
\item 2026-10-30 (P2-T002): Proposed sealed statements \texttt{P2.Restoration}, \texttt{P2.Afforestation} and \texttt{P2.Harmonization},
with rationales, sanity examples and negative controls. Added a method section for writing statements.
Condensed the Carbon Continuity material into its own section, and retitled the paper to cover the P2 program.
\end{itemize}

\bibliographystyle{plainnat}
\bibliography{refs}
\end{document}
